% GENERATED FILE. Edit canonical lessons, then run npm run generate:books.
% canonical-source-hash: fnv1a64:3f2d67e9e82b83f6
% canonical-lessons: AR-C326-silah, AR-C326-sayf, AR-C326-rumh, AR-C326-sahm, AR-C326-dir

\chapter{Old weapons}
\label{ch:ar-old-weapons}
\hlchaptermodality{326}

\begin{chapteropening}
\textbf{By the end of this chapter:} \emph{I can say a weapon, a sword, a spear, an arrow and armour.}

Complete the last lesson of chapter 326: I can say a weapon, a sword, a spear, an arrow and armour.
\end{chapteropening}

\section[\texorpdfstring{\ar{سلاح} (silāḥ)}{سلاح (silāḥ)}]{\ar{سلاح} (silāḥ) — a weapon}

\label{lesson:AR-C326-silah}



\begin{quote}
Before the new one: say the Arabic for to confess, then the Arabic for to deny.
\end{quote}

\subsection*{You'll want to know: \ar{سلاح}}
\textbf{\ar{سلاح}} — \emph{silāḥ} — \textquotedblleft{}a weapon\textquotedblright{}.

\begin{quote}
\ar{كان} \ar{مع} \ar{الحارس} \ar{سلاح} \ar{صغير}.
\end{quote}

\begin{quote}
\emph{kāna maʿa al-ḥāris silāḥ ṣaghīr.}
\end{quote}

\begin{quote}
The guard had a small weapon.
\end{quote}

Plural: \ar{أسلحة}, \emph{asliḥa}.

\begin{cousinweb}
\textbf{\ar{سلاح}} (\emph{silāḥ}) is built on the native Arabic root \emph{s-l-ḥ} (Lane), where \emph{silāḥ} names the arms of war and \emph{musallaḥ}, armed, is its participle.
\end{cousinweb}

\begin{grammarlens}[title={What you've built}]
One more word for telling what happened.
\end{grammarlens}

\subsection*{Your turn}
\begin{itemize}
\raggedright
  \item \emph{Say it:} \emph{silāḥ}
  \item \emph{Say it:} \emph{silāḥ}, once more
  \item \emph{Say it:} \emph{kāna maʿa al-ḥāris silāḥ ṣaghīr.} — The guard had a small weapon.
  \item \emph{Recall:} say the Arabic for to confess, then the Arabic for to deny, then say \emph{silāḥ} again
\end{itemize}

\begin{tcolorbox}[breakable,colback=teal!4,colframe=teal!35!black,title={Before you move on}]
What does \ar{سلاح} mean? (\textquotedblleft{}A weapon\textquotedblright{}.) Say it once more.
\end{tcolorbox}
\section[\texorpdfstring{\ar{سيف} (sayf)}{سيف (sayf)}]{\ar{سيف} (sayf) — a sword}

\label{lesson:AR-C326-sayf}



\begin{quote}
Before the new one: say the Arabic for to deny, then the Arabic for a weapon.
\end{quote}

\subsection*{You'll want to know: \ar{سيف}}
\textbf{\ar{سيف}} — \emph{sayf} — \textquotedblleft{}a sword\textquotedblright{}.

\begin{quote}
\ar{علق} \ar{الفارس} \ar{سيفه} \ar{على} \ar{الجدار}.
\end{quote}

\begin{quote}
\emph{ʿallaqa al-fāris sayfahu ʿalā al-jidār.}
\end{quote}

\begin{quote}
The knight hung his sword on the wall.
\end{quote}

Plural: \ar{سيوف}, \emph{suyūf}.

\begin{cousinweb}
\textbf{\ar{سيف}} (\emph{sayf}) is filed by Lane under the root \emph{s-y-f}. A link with Greek \emph{xiphos}, a sword, is often proposed, but the direction of the borrowing, if any, is not settled.
\end{cousinweb}

\begin{grammarlens}[title={What you've built}]
One more word for telling what happened.
\end{grammarlens}

\subsection*{Your turn}
\begin{itemize}
\raggedright
  \item \emph{Say it:} \emph{sayf}
  \item \emph{Say it:} \emph{sayf}, once more
  \item \emph{Say it:} \emph{ʿallaqa al-fāris sayfahu ʿalā al-jidār.} — The knight hung his sword on the wall.
  \item \emph{Recall:} say the Arabic for to deny, then the Arabic for a weapon, then say \emph{sayf} again
\end{itemize}

\begin{tcolorbox}[breakable,colback=teal!4,colframe=teal!35!black,title={Before you move on}]
What does \ar{سيف} mean? (\textquotedblleft{}A sword\textquotedblright{}.) Say it once more.
\end{tcolorbox}
\section[\texorpdfstring{\ar{رمح} (rumḥ)}{رمح (rumḥ)}]{\ar{رمح} (rumḥ) — a spear, a lance}

\label{lesson:AR-C326-rumh}



\begin{quote}
Before the new one: say the Arabic for a weapon, then the Arabic for a sword.
\end{quote}

\subsection*{You'll want to know: \ar{رمح}}
\textbf{\ar{رمح}} — \emph{rumḥ} — \textquotedblleft{}a spear, a lance\textquotedblright{}.

\begin{quote}
\ar{حمل} \ar{الصياد} \ar{رمحا} \ar{طويلا}.
\end{quote}

\begin{quote}
\emph{ḥamala aṣ-ṣayyād rumḥan ṭawīlan.}
\end{quote}

\begin{quote}
The hunter carried a long spear.
\end{quote}

Plural: \ar{رماح}, \emph{rimāḥ}.

\begin{cousinweb}
\textbf{\ar{رمح}} (\emph{rumḥ}) is built on the native Arabic root \emph{r-m-ḥ} (Lane), where \emph{ramaḥa} is \textquotedblleft{}to pierce with a spear\textquotedblright{} and is said of a horse that kicks.
\end{cousinweb}

\begin{grammarlens}[title={What you've built}]
One more word for telling what happened.
\end{grammarlens}

\subsection*{Your turn}
\begin{itemize}
\raggedright
  \item \emph{Say it:} \emph{rumḥ}
  \item \emph{Say it:} \emph{rumḥ}, once more
  \item \emph{Say it:} \emph{ḥamala aṣ-ṣayyād rumḥan ṭawīlan.} — The hunter carried a long spear.
  \item \emph{Recall:} say the Arabic for a weapon, then the Arabic for a sword, then say \emph{rumḥ} again
\end{itemize}

\begin{tcolorbox}[breakable,colback=teal!4,colframe=teal!35!black,title={Before you move on}]
What does \ar{رمح} mean? (\textquotedblleft{}A spear, a lance\textquotedblright{}.) Say it once more.
\end{tcolorbox}
\section[\texorpdfstring{\ar{سهم} (sahm)}{سهم (sahm)}]{\ar{سهم} (sahm) — an arrow; a share}

\label{lesson:AR-C326-sahm}



\begin{quote}
Before the new one: say the Arabic for a sword, then the Arabic for a spear.
\end{quote}

\subsection*{You'll want to know: \ar{سهم}}
\textbf{\ar{سهم}} — \emph{sahm} — \textquotedblleft{}an arrow; a share\textquotedblright{}.

\begin{quote}
\ar{رمى} \ar{الصياد} \ar{سهما} \ar{نحو} \ar{الشجرة}.
\end{quote}

\begin{quote}
\emph{ramā aṣ-ṣayyād sahman naḥwa ash-shajara.}
\end{quote}

\begin{quote}
The hunter shot an arrow at the tree.
\end{quote}

Plural: \ar{سهام}, \emph{sihām}.

In business a share is \ar{سهم}, plural \ar{أسهم}. A road sign arrow is \ar{سهم} too.

\begin{cousinweb}
\textbf{\ar{سهم}} (\emph{sahm}) is built on the native Arabic root \emph{s-h-m} (Lane), where \emph{sahm} is an arrow, and an arrow drawn as a lot: so also a share, \emph{sāhama}, to take part.
\end{cousinweb}

\begin{grammarlens}[title={What you've built}]
One more word for telling what happened.
\end{grammarlens}

\subsection*{Your turn}
\begin{itemize}
\raggedright
  \item \emph{Say it:} \emph{sahm}
  \item \emph{Say it:} \emph{sahm}, once more
  \item \emph{Say it:} \emph{ramā aṣ-ṣayyād sahman naḥwa ash-shajara.} — The hunter shot an arrow at the tree.
  \item \emph{Recall:} say the Arabic for a sword, then the Arabic for a spear, then say \emph{sahm} again
\end{itemize}

\begin{tcolorbox}[breakable,colback=teal!4,colframe=teal!35!black,title={Before you move on}]
What does \ar{سهم} mean? (\textquotedblleft{}An arrow; a share\textquotedblright{}.) Say it once more.
\end{tcolorbox}
\section[\texorpdfstring{\ar{درع} (dirʿ)}{درع (dirʿ)}]{\ar{درع} (dirʿ) — armour; a shield}

\label{lesson:AR-C326-dir}



\begin{quote}
Before the new one: say the Arabic for a spear, then the Arabic for an arrow.
\end{quote}

\subsection*{You'll want to know: \ar{درع}}
\textbf{\ar{درع}} — \emph{dirʿ} — \textquotedblleft{}armour; a shield\textquotedblright{}.

\begin{quote}
\ar{لبس} \ar{الجندي} \ar{درعا} \ar{من} \ar{الحديد}.
\end{quote}

\begin{quote}
\emph{labisa al-jundī dirʿan min al-ḥadīd.}
\end{quote}

\begin{quote}
The soldier wore armour of iron.
\end{quote}

Plural: \ar{دروع}, \emph{durūʿ}.

Today a \ar{درع} is also a shield-shaped trophy.

\begin{cousinweb}
\textbf{\ar{درع}} (\emph{dirʿ}) is built on the native Arabic root \emph{d-r-ʿ} (Lane), where \emph{dirʿ} is a coat of mail, rings linked over the body.
\end{cousinweb}

\begin{grammarlens}[title={What you've built}]
One more word for telling what happened.
\end{grammarlens}

\subsection*{Your turn}
\begin{itemize}
\raggedright
  \item \emph{Say it:} \emph{dirʿ}
  \item \emph{Say it:} \emph{dirʿ}, once more
  \item \emph{Say it:} \emph{labisa al-jundī dirʿan min al-ḥadīd.} — The soldier wore armour of iron.
  \item \emph{Recall:} say the Arabic for a spear, then the Arabic for an arrow, then say \emph{dirʿ} again
\end{itemize}

\begin{tcolorbox}[breakable,colback=teal!4,colframe=teal!35!black,title={Before you move on}]
What does \ar{درع} mean? (\textquotedblleft{}Armour; a shield\textquotedblright{}.) Say it once more.
\end{tcolorbox}
